\section{Discussion and Conclusion}\label{sec:conclusion}

Across 100+ experiments and 6000+ evaluation episodes on two VLA model families, three findings converge: CAR-based VLA safety evaluation on SafeLIBERO is structurally unreliable. Action collapse renders reported CAR illusory (\S\ref{sec:dfcar}); EE-only oracle refinement is counterproductive (\S\ref{sec:ee_only}); and collision-free demonstrations account for nearly all measured avoidance gains, with EE-only SDF contributing no statistically significant improvement (\S\ref{sec:avoidance_source}).

\paragraph{Why EE-only SDF fails.}
Two independent causes explain the null result. \textit{(i)~Demonstration data subsumes the SDF signal.} Collision-free demonstrations implicitly encode spatial avoidance; the ${\sim}$26\,pp CAR gain arises from this distributional shift, with per-task ordering invariant across all SDF conditions (\S\ref{sec:avoidance_source}). Even at $\lambda = 0.1$, CAR remains within the null distribution ($p = 0.131$). \textit{(ii)~EE monitoring cannot prevent full-body collisions.} The SDF gradient acts on a single point while collisions occur along intermediate links; ground-truth oracle refinement \emph{reduces} CAR by 5--10\,pp (\S\ref{sec:ee_only}) because pushing the EE away rotates unmonitored links through obstacles. These causes are independent; effective geometric safety requires both full-body supervision and training signals beyond demonstrations. Constrained-optimization approaches (e.g., SafeVLA) address this through a fundamentally different paradigm---directly constraining violation probability---while runtime filtering sidesteps training altogether.

\paragraph{Implications for evaluation.}
dfCAR corrects action-collapse contamination using only EE position logging, already available in most evaluation pipelines. The bimodal displacement distribution makes $\tau_d$ robust across five orders of magnitude. We recommend reporting both CAR and dfCAR; the gap quantifies contamination across a ${\sim}$4000$\times$ parameter range (Pi0 Full-FT ${\sim}$2.5B to Pi0+A1 ${\sim}$658K trainable). More broadly, safety may need to operate at the planning or joint-space constraint layer, since any EE-centric loss inherits the same blind spot.

\paragraph{Limitations and scope.}
All experiments use MuJoCo simulation (SafeLIBERO, two VLA families, four pick-and-place tasks). The action-collapse mechanism is structurally general---any policy prone to action collapse will inflate CAR---but its prevalence across other benchmarks and real hardware remains untested. The EE-only SDF null result is specific to auxiliary loss integration and should not be extrapolated to alternative pathways (Section~\ref{sec:related}). The sim-to-real gap introduces uncertainty: collision detection precision differs on physical hardware and sensor noise may affect dfCAR differently. Statistical power is limited by $n{=}3$ seeds per condition (MDE ${\approx}$ 5\,pp at 80\% power). Oracle experiments used largely action-collapsed models; validation on actively moving policies remains future work. Full-body SDF remains untested.

\paragraph{Future work.}
Three directions follow: (1) full-body SDF supervision along the kinematic chain; (2) cross-architecture evaluation on OpenVLA, RT-2, and other VLA families; (3) planning-layer safety constraints in joint space rather than EE space.

\paragraph{Broader impact.}
dfCAR is a diagnostic complement to standard CAR, not a replacement; practitioners should report both to distinguish genuine avoidance from action-collapse artifacts. Our finding that EE-only SDF is redundant at the tested scale is a methodology insight, not a verdict against safety-aware training---geometric supervision may prove effective through alternative pathways or full-body collision surfaces.
